Semiconductor materials capable of complementing Si in nanoelectronics are needed to face the recent global chip manufacturing crisis. Significant efforts have been expended in adapting graphene electronics into integrated circuits, but the various limitations inherent to
its operating mechanism incentivise the study of other carbon nanostructures.

Here, I investigate molecular graphene nanostructures as potential candidate to this end.
I began by describing the instrumentation and device designs used to collect the electrical transport data presented in this thesis. I then demonstrate that molecular graphene nanoribbons (MGNRs) can be employed as a molecular quantum dot platform at high temperatures.
These findings provide the first observation of quantum events at high temperatures in a full-carbon, all-around-molecular structure at the device level.

Compared to carbon nanotubes, promoting MGNRs de-bundling when they are dispersed in organic solvent can be achieved more easily, and I describe possible strategies that could optimize the device yield and reproducibility.

I then examine the effect of adding dopant groups to the ribbon edges. The energy shift induced by the dopant to the ribbon Fermi level is, however, weak and useful at the device level. I present the first field-effect transistor measurements of an unique class of materials in which the dopant atoms are inserted directly into the honeycomb lattice, and propose possible improvements.

Based on the research I conducted for this thesis, I believe that the future of practical MGNR-based electronics is still in its infancy. They host a multitude of physical events that should be investigated thoroughly using both STM and device-level techniques. Improving the length of the ribbon while maintaining a high degree of debundling will be crucial to the development of the first MGNR-based microprocessor.

On the basis of the research I performed in this thesis, I believe that the outlook for any practical MGNR-based electronics is distant. MGNRs host a plethora of physics phenomena which they should be first investigate thouroghly, both through STM and at the device level. The necessity of synthesising longer ribbon while mantaining a high degree of debundling will also be fundamental towards the realisation of a first MGNR-based microprocessor.